\chapter{Advanced types}\label{ch-advanced}

The types of the previous chapters were built in simple ways: arrays of
numbers, records of strings and numbers, functions of numbers. This
chapter shows the full power of the type system. It introduces arrays
that change their size, unions of types, references to references, and
recursive types such as lists and trees, and it combines them in larger
examples. It then collects the rules that earlier chapters introduced
piece by piece: all implicit conversions and the places where each is
allowed, and the rules that govern how long a value lives.

This is the chapter where Algol 68 differs most from the languages that
most programmers know. Many of its rules are easy to state and have
consequences that are not obvious, so the chapter shows those
consequences, including the ones that the compilers do not report.

The guiding principle behind all of it is \term{orthogonality}
(Section~\ref{sec-orthogonality}): every way of building a type can be
applied to every type, and every construct works the same way whatever
the types involved.

\section{Resizable arrays}\label{sec-flex}

The arrays of Chapter~\ref{ch-data} have fixed bounds. A variable declared with
\code|flex| before the bounds can hold arrays of any size; each assignment
gives it the bounds of the assigned array:

\example{ch07/flex}

\code|flex [1:0] int list| declares a variable that holds an array of
\code|int|, initially empty. The bounds in the declaration only determine
the initial array. The Report calls such a variable a \rr{flexible name}.

The value of a resizable array is an ordinary array: \code|snapshot| is
a copy with the bounds 1 to 2, which does not change when \code|list|
changes. The type of such a value is the array type without
\code|flex|; the Report says that the type \rr{deflexes} to it.

There is no operation that adds an element to a resizable array in
place. An array literal cannot help either: each element of the literal
must be an \code|int|, and \code|squares| is an array, so the compiler
rejects \code|(squares, i * i)|:

\example{ch07/append}

Instead, a new, longer array is built and assigned. A procedure does this
once and for all:

\example{ch07/append-proc}

The parameter has the type \code|ref flex [] int|: a reference to a
resizable array variable, the only kind of variable whose bounds may
change. A \code|ref [] int| parameter would accept fixed arrays as well,
and could therefore not be given a longer array. In the Report, the two
are even different types: a \code|ref flex [] int| cannot be passed as a
\code|ref [] int|. This is the price of Algol 68's value semantics for
arrays: what other languages call a dynamic array or list
(\texttt{std::vector}, \texttt{ArrayList}, a Python \texttt{list}) is
here written out explicitly.

\begin{implnote}
ga68, a68j and Algol 68 Genie accept a \code|ref flex [] int| argument for
a \code|ref [] int| parameter, contrary to the Report. A program that
relies on this is not portable.
\end{implnote}

\section{Where \textbf{flex} can stand}\label{sec-flex-where}

\code|flex| is a property of variables, not of values: an array value
always has fixed bounds, and only a variable can be given a new array of
a different size. \code|flex| therefore belongs in the type of a
variable, that is, under a \code|ref|. It can appear in a variable
declaration, as above, in the type \code|ref flex [] int| of a
parameter, and in a type declaration that is used to declare variables:
\begin{lstlisting}
mode Buffer = flex [1:0] int;
Buffer b := (1, 2, 3);
b := (4, 5)
\end{lstlisting}
The standard type \code|string| is exactly such a declaration
(Section~\ref{sec-string}). A type declaration needs bounds, even if they
are only initial ones: \code|mode B = flex [] int| is rejected with
``actual bounds expected''. A constant cannot be resizable, so
\code|flex [] int b = (1, 2, 3)| is not allowed by the Report; the
implementations accept it and treat \code|b| as a \code|[] int|.

Deflexing also applies to the fields of a record. A variable of type
\code|struct (string name, int age)| holds records whose \code|name| can
have any length, but the record value that the variable holds has a
field of type \code|[] char|, with the bounds of its current name.

\subsection*{Jagged arrays}

A two-dimensional array \code|[,] int| is rectangular: all its rows have
the same length. An array of arrays, \code|[] [] int|, is a different
type, whose elements are arrays in their own right and may have
different lengths, like a Java \code|int[][]|. With \code|flex| at both
levels, both the number of rows and each row can change:

\example{ch07/jagged}

In the literal \code|((1, 2), 3, (4, 5, 6))|, the element \code|3| is
converted to an array with one element (rowing,
Section~\ref{sec-rowing}). The assignment to \code|rag[2]| gives the second
row a new length, and \code|rag := rag[1:2]| drops the last row. In
between lie the type \code|[3] flex [1:0] int|, three rows of which each
can change, and \code|flex [1:0] [3] int|, a varying number of rows of
fixed length 3.

\begin{implnote}
ga68 stops with an internal compiler error when a value of type
\code|[] [] int| that is not a literal is assigned to a variable of type
\code|flex [] flex [] int|, as in \code|rag := jag|. Assigning a literal,
as in the example, works.
\end{implnote}

\section{Growing an array efficiently}\label{sec-vector}

The procedure \code|append| of Section~\ref{sec-flex} copies the whole
array each time it adds an element. Building an array of $n$ elements
this way copies about $n^2/2$ elements, which is slow for large $n$. The
dynamic arrays of other languages avoid this by keeping some spare room:
the array is larger than the number of elements in use, and when it is
full, it is replaced by one twice as large. Each element is then copied
only a few times on average.

In Algol 68 the same structure is a record with a resizable array and a
count:

\program{ch07/vector}

\noindent The program prints the number of elements and the capacity
after each \code|push|:
\outputof{ch07/vector}

The field \code|data| of a \code|Vector| variable is itself a resizable
array variable: if \code|v| has the type \code|ref Vector|, then
\code|data of v| has the type \code|ref flex [] int| and can be assigned
to, indexed, or passed to \code|append| of Section~\ref{sec-flex}. The
parentheses in \code|(data of v)[i]| are needed because indexing binds
more tightly than selection (Section~\ref{sec-multiple-selection}).

The same cost applies to strings: \code|s +:= c| in a loop builds a new
string each time. For long texts, collecting the characters in a buffer
of this kind and converting them at the end is much faster.

\section{Strings}\label{sec-string}

The type \code|string| is defined in the standard library as
\code|flex [1:0] char|, a resizable array of characters. That is why a
\code|string| variable can hold strings of any length. The library adds a
few operators:

\example{ch07/strings}

\code|+| joins strings and characters; \code|s +:= t| appends \code|t|
to \code|s|, and \code|t +=: s| puts \code|t| in front of \code|s|.
\code|*| repeats a string a given number of times, and \code|s *:= n|
replaces \code|s| by \code|n| copies of itself. The comparisons
\code|=|, \code|/=|, \code|<|, \code|<=|, \code|>|, \code|>=| compare
strings in dictionary order, by the codes of their characters. They are
defined for \code|[] char|, so they apply to every array of characters,
not only to \code|string| variables.

\section{References into resizable arrays}\label{sec-transient}

An element of an array variable is itself a variable: \code|a[1]| is a
reference to the first element. For a resizable array, such a reference
must not be kept, since the next assignment to the array may give it a
different size and place the elements elsewhere; the reference would then
refer to a cell that no longer belongs to the array. The language
therefore forbids binding such a reference to a name or storing it:

\example{ch07/transient}

The Report calls a reference into a resizable array a \rr{transient
name}. The same problem is known in C++ as iterator invalidation; Algol
68 rules it out at compile time.

The rule is stricter than it looks. A transient reference may be used
only directly, for an assignment or to obtain its value. It may not be
bound to a name, stored, returned, or passed to a procedure or an
operator, since a parameter is bound like a constant declaration and no
type can be written for a transient reference. That excludes some
everyday code:

\example{ch07/transient-count}

\code|+:=| is an operator whose left operand is a reference, so
\code|count[3] +:= 1| passes a transient reference to it. For the same
reason \code|read (count[3])| and \code|swap (count[i], count[j])| are
rejected, and so is \code|s[i] +:= "x"| for an element of a
\code|flex [] string|. Note that \code|string| is a resizable array, so
the rule applies to the characters of every \code|string| variable.
Selecting a field across all elements, \code|x of path| for a
\code|flex [] Point| variable \code|path|, gives a transient reference
as well.

There are three ways around it. The plain assignment
\code|count[3] := count[3] + 1| is allowed. If the size is known once the
array has been filled, a fixed array can take over:
\code|[UPB count] int c := count| is an ordinary array whose elements may
be passed anywhere. And the restriction applies only to resizable
arrays: the elements of fixed arrays and the fields of records are
ordinary variables.

\section{Sharing arrays}\label{sec-sharing}

Arrays are values in Algol 68: assigning an array copies its elements
(Section~\ref{sec-array-values}). An array is shared only when a
reference to it is shared explicitly, which the type always shows:

\program{ch07/share}

\noindent prints
\outputof{ch07/share}

\code|r| is a constant of type \code|ref [] int| bound to the variable
\code|a|: another name for the same array, so \code|r[1] := 10| changes
\code|a|. \code|s| is a variable that holds a reference to a part of
\code|a|; \code|s[1]| is \code|a[2]|. And \code|one| shows the rowing of
a variable: the reference \code|x| becomes a reference to an array with
one element, which is \code|x| itself.

The same distinction exists for single values. \code|ref int r = x|
makes \code|r| another name for \code|x|, like a reference
\code|int& r = x| in C++. \code|ref int p := x| creates a variable that
holds a reference to \code|x| and can later refer to something else, like
a pointer \code|int* p = &x| in C. In Java and Python, objects and lists
are always shared and a copy must be made explicitly; in Algol 68, as in
C++ (and, with moves instead of copies, in Rust), values are copied and
sharing must be made explicit.

\section{Unions}\label{sec-unions}

A \term{union} type holds a value of one of several types, together with
an indication of which type it is. \code|union (int, real, string)| holds
an \code|int|, a \code|real| or a \code|string|. A value of any of these
types can be assigned to a union variable or passed as a union parameter;
this implicit conversion is \term{uniting} (\rr{uniting}).

To use the value inside a union, the program must find out its type. This
is done with a special form of \code|case|, which the Report calls a
\rr{conformity clause}:

\example{ch07/union}

Each alternative names a type and, optionally, a name for the value:
\code|(int n): ...| is chosen if the union holds an \code|int|, and binds
\code|n| to it. An \code|out| part handles the types that no alternative
names. A union without a name in its alternative only tests the type:

\example{ch07/union-test}

The union is a \term{tagged union}, or \term{sum type}: the type of the
current value is recorded with it and checked at run time, so a value can
never be read as the wrong type. It corresponds to the \texttt{enum} of
Rust and Swift, the variant types of ML and Haskell and
\texttt{std::variant} of C++, and the conformity clause corresponds to
their pattern matching. Unlike those, an Algol 68 union has no names for
its alternatives; the alternatives are distinguished by their types only.
Two alternatives that carry the same type of value must therefore be
wrapped in different record types.

The order of the types in a union does not matter, and a union that
contains a union contains the members of the inner one:
\code|union (int, real)| and \code|union (real, int)| are the same type,
and \code|union (int, union (real, string))| is
\code|union (int, real, string)|. A union value can be converted
implicitly to a larger union, but not to a smaller one, since the value
might be of a type that the smaller union lacks.

The types of a union must be clearly distinct: \code|union (int, ref int)|
is rejected (``firmly related components''), because an \code|int|
variable could become either member. A union value needs the space of
its largest member plus the tag, and every conformity clause is a test at
run time.

The \code|print| procedure is an example of unions at work: its parameter
is an array of a union of all the types it can print.

\section{Working with unions}\label{sec-union-details}

\subsection*{What a conformity clause binds}

The name in an alternative is a constant that holds a copy of the value:
in \code|(int n): n +:= 1| the compiler rejects the assignment. To
change the value in a union variable \code|u|, assign a new value to
\code|u| itself, as in \code|(int n): u := n + 1|. An alternative may also
name a union of some of the members: \code|(union (int, real) x): ...|
is chosen for an \code|int| or a \code|real|, and \code|x| is then a
value of that smaller union.

The alternatives are chosen by the exact type of the value, without
conversions. An \code|int| in a union does not match an alternative for
\code|real|:

\example{ch07/nowiden}

\begin{pitfall}
A conformity clause without an \code|out| part, none of whose
alternatives matches, does nothing and yields an undefined value. There
is no error at compile time or at run time:
\begin{lstlisting}
union (int, real, string) u := "seven";
int n = case u in (int i): i, (real r): ROUND r esac
\end{lstlisting}
ga68 and a68j give \code|n| the value 0, Algol 68 Genie the value 1.
Rust rejects a match that does not cover every case, and a Haskell
compiler warns about it and the program stops at run time; Algol 68 does
neither. Every conformity clause that is not meant to ignore
the remaining types should have an \code|out| part that reports the
error.
\end{pitfall}

\subsection*{Optional values}

A union with \code|void| is the Algol 68 form of an optional value, the
\code|Option<T>| of Rust or the \code|Optional<T>| of Java. The only value
of type \code|void| is written \code|empty|:

\example{ch07/maybe}

\code|nil| serves the same purpose for references, but only for them; a
search that returns an index needs \code|union (int, void)|. The
\code|while| part of the loop uses the short form of the conformity
clause: \code|(result | (void): true | false)| is
\code|case result in (void): true out false esac|. Rosetta Code
programs use the short form often.

\subsection*{Which values can be united}

Implicit conversions happen in a fixed order (Section~\ref{sec-contexts}):
uniting and widening are alternatives, and rowing can only follow them.
A value is therefore united only if its own type is a member of the
union, not a type that it could be widened or rowed to:

\example{ch07/union-order}

\example{ch07/union-char}

\code|1| is an \code|int|, not a \code|real|, and \code|"x"| is a
\code|char|, not a \code|string|; neither is a member of the union. The
remedies are to write the value with the right type (\code|1.0|) or to
annotate it (\code|real (1)|, \code|string ("x")|,
Section~\ref{sec-casts}). For the same reason, \code|print (1)| prints an
integer: \code|int| is a member of the union that \code|print| accepts,
so it is not widened. A one-character string literal must be annotated
when it is to become a \code|string| in a union; this is a common
surprise.

A record literal has no type of its own, so it cannot be united either:
with \code|mode Pt = struct (int x, y), Sz = struct (int w, h)|, the
declaration \code|union (Pt, Sz) r := (1, 2)| is rejected, and
\code|union (Pt, Sz) r := Pt (1, 2)| says which member is meant.
Uniting does not reach through a reference: a function with a
parameter of type \code|ref union (int, real)| does not accept an
\code|int| variable; the caller must pass a variable of the union type. Nor can a union whose
members are arrays be indexed directly: a conformity clause must first
say which member it holds. And no
conversion applies to the elements of an array: a \code|[] int| is never
converted to a \code|[] real| or to an array of a union. (Java allows an
\code|Integer[]| to be used as an \code|Object[]|, which is why Java must
check every store into an array at run time.)

\section{References to references}\label{sec-pointers}

A reference is a value, so there are variables that hold references. A
variable of type \code|ref int| is declared like any other variable, and
its type is \code|ref ref int|. Such a variable is what C calls a
pointer:

\example{ch07/pointers}

The assignment \code|p := b| stores the reference \code|b| in \code|p|;
it does not change \code|a| or \code|b|. Which cell an assignment changes
is determined by the type of its left-hand side: the left-hand side of
\code|p := b| is \code|p| itself, of type \code|ref ref int|, so a
\code|ref int| is stored. To assign to the cell that \code|p| refers to,
the left-hand side must be the \code|ref int| that \code|p| holds. The
expression \code|ref int (p)| says exactly that: it is \code|p| converted
to the type \code|ref int|, which dereferences it once.

The \term{null reference} is written \code|nil|; it refers to no cell, and
any attempt to use it as a variable is an error. Two references are
compared with \code|is| and \code|isnt| (also written \code|:=:| and
\code|:/=:|); the Report calls these \rr{identity relations}. They
compare references, not the values in the cells: \code|p is b| is true
because \code|p| refers to the same cell as \code|b|.

Comparing a pointer with \code|nil| needs care. In an identity relation,
one side is in a soft context, where no dereferencing happens, and the
other in a strong one, in either order. \code|nil| can only stand in a
strong context, so in \code|q is nil| the variable \code|q| is the soft
side and is not dereferenced at all: the comparison is between the
variable \code|q| itself and \code|nil|, which is always false. The test
for a null pointer is \code|ref int (q) is nil|:

\example{ch07/nil}

This is the most common pitfall of Algol 68 references, and a direct
consequence of the rule that an implicit conversion never happens
unless the context requires it.

\section{Comparing references}\label{sec-identity}

The same rule applies to two pointer variables. If \code|p| and \code|q|
are both of type \code|ref ref int| and both refer to \code|n|, then
\code|p is q| compares the two variables themselves, which are different
cells:

\example{ch07/refcompare}

\noindent prints \code|F T T|. The question that is usually meant, whether
the two pointers refer to the same cell, is
\code|ref int (p) is ref int (q)|. The rule of thumb: annotate one side
with the type of the references that are to be compared, and the
other side is dereferenced to the same type. \code|ref int (p) is n|
compares the reference held by \code|p| with the variable \code|n|.

An identity relation is not an operator. It cannot be redefined, and it
binds more loosely than every operator, so it must be put in parentheses
when it is combined with \code|AND|, \code|OR|, \code|andth| or
\code|orel|: \code|p is q AND b| means \code|p is (q AND b)| and is
rejected, and the condition of a loop over a list is written
\code|while (ref Node (p) isnt no_node) andth value of p /= v do|.

\section{Records with reference fields}\label{sec-ref-fields}

A record field can be a reference, and a constant record can then
contain something that changes. A gas meter has a fixed serial number and
a reading that grows:

\example{ch07/meter}

\noindent prints \code|30|. \code|m| is a constant, but its field
\code|reading| is a \code|ref int|, and the cell it refers to can be
assigned to. Copying the record into \code|v| copies the reference, not
the cell, so \code|m| and \code|v| share the reading. This is the
behavior of a \code|final| field that refers to a mutable object in
Java. The semaphores of Chapter~\ref{ch-par} are records of this kind.

Assigning to the reading through the variable \code|v| needs an
annotation. \code|reading of v| is a field of a variable, and therefore
has the type \code|ref ref int|: the variable that holds the reference.
\code|reading of v := 30| would try to store an \code|int| in it. The
annotation \code|ref int (reading of v)| dereferences once and yields
the cell to be changed. In contrast, selecting from or indexing through
a reference needs no annotation: if \code|pa| is a variable of type
\code|ref [] int|, then \code|pa[1] := 5| assigns to an element of the
array it refers to, and \code|next of p| in the lists below selects from
the record that \code|p| refers to. The Report calls this the weak
dereferencing of Section~\ref{sec-contexts}.

\section{Type annotations}\label{sec-casts}

An expression of the form \code|Type (expression)| is a \term{cast}. The
word suggests a conversion as in C, but an Algol 68 cast converts nothing
by force: it only places the expression in a strong context with the
given target type, so that exactly the implicit conversions to that type
happen. A cast is therefore a type annotation, like \texttt{e: T} in
Rust or TypeScript. It cannot convert a \code|real| to an \code|int|; that
needs \code|ROUND| or \code|ENTIER|.

\example{ch07/casts}

The cast \code|real (1)| makes the united value a \code|real| rather than
an \code|int|. The declaration \code|proc int saved = next| binds the
function itself, since its type is a function type, whereas
\code|int value = next| calls it.

An array or record literal is allowed only where the type it must have
is known, that is, in a strong context. It cannot be the operand of an
operator, where the operator is not yet known; an annotation provides the
type:

\example{ch07/display-operand}

\code|UPB (1, 2, 3)| would be rejected. Annotations are also the way to
say which member of a union a record literal belongs to
(Section~\ref{sec-union-details}), and \code|void (x := f)| discards a
value explicitly.

\section{Orthogonality}\label{sec-orthogonality}

The constructions of types can be combined without restriction. There
are arrays of functions (\code|[] proc (real) real|), functions that
return functions (\code|proc (int) proc (int) int|), records with
function fields, unions of records, references to arrays of references,
and so on. Every type except \code|void| can be the type of a parameter,
a result, a field or an element; \code|void| can be the result of a
function and a member of a union.

The constructs of the language are equally orthogonal. Any expression
can stand wherever a value of its type is required, and a block or an
\code|if| can stand where other languages only allow a simple
expression, as Chapter~\ref{ch-control} showed. This is why Algol 68 needs few special
rules, and why its few rules apply everywhere.

\section{Implicit conversions and contexts}\label{sec-contexts}

The chapters so far introduced the six implicit conversions one at a time.
Here they are together:

\begin{center}
\begin{tabular}{lll}
\toprule
Conversion & Report & Effect \\
\midrule
dereferencing & \rr{dereferencing} & a reference is replaced by the value it refers to \\
calling & \rr{deproceduring} & a parameterless function is called \\
uniting & \rr{uniting} & a value becomes a value of a union type \\
widening & \rr{widening} & \code|int| to \code|real| to \code|compl|; \\
         &              & \code|bits| to \code|[] bool|; \code|bytes| to \code|[] char| \\
rowing & \rr{rowing} & a value becomes an array with one element, \\
       &             & an array gains a dimension \\
discarding & \rr{voiding} & a value is thrown away \\
\bottomrule
\end{tabular}
\end{center}

Rowing also applies to references: a variable of type \code|ref int|
can become a \code|ref [] int| whose only element is the variable itself
(Section~\ref{sec-sharing}).

Where each conversion is allowed depends on the context in which the
expression stands. There are five contexts, from the most to the least
permissive:

\begin{center}
\begin{tabular}{lll}
\toprule
Context & Where & Conversions allowed \\
\midrule
strong & right-hand side of a declaration or & all \\
       & assignment, argument, cast, function & \\
       & body, element of a literal, one & \\
       & side of an identity relation & \\
firm   & operand of an operator & dereferencing, calling, uniting \\
meek   & condition of \code|if|, \code|case|, \code|while|; & dereferencing, calling \\
       & index, bound, \code|from|, \code|by|, \code|to|; & \\
       & the function in a call & \\
weak   & the array in an indexing, the record & dereferencing (but not of \\
       & in a selection & the last reference), calling \\
soft   & left-hand side of an assignment, & calling \\
       & one side of an identity relation & \\
\bottomrule
\end{tabular}
\end{center}

In a strong context, the required type is known, and the compiler applies
whatever conversions lead to it, always in the same order: first
dereferencing and calling as often as needed, then either one uniting or
widening (twice for \code|int| to \code|compl|), then rowing (once per
dimension added). A parameterless function \code|f| of type
\code|proc ref int| that stands where a \code|real| is required is
called, its result is dereferenced, and the \code|int| is widened.

In a firm context, the operator is not
known until the types of the operands are, so the conversions that would
change the type substantially (widening, rowing) are excluded; uniting is
allowed so that an operator can be declared for a union. A meek context
requires a value of a particular kind, such as a \code|bool| for a
condition or an \code|int| for an index, and only dereferencing and
calling are allowed. A weak context keeps a reference: in \code|a[i] := 0|,
the array \code|a| is not dereferenced to its value, since the element
must remain a reference to be assigned to; this is also why a slice of an
array variable is a reference (Section~\ref{sec-slices}). And a soft
context keeps the reference on the left-hand side of an assignment
entirely.

The branches of an \code|if| or a \code|case| that yields a value are
balanced: one of them is in the context of the whole expression, and the
others are in a strong context, so they can be converted to a common
type even where the context itself allows little. \code|if b then 1 else 2.5 fi|
is a \code|real|, and so is \code|(b | 1 | 2.5) + 1|. On the left-hand
side of an assignment, \code|if b then x else y fi := 0| needs branches
that yield references of one type, though a branch may be dereferenced to
match the other.

The Report calls a context a \rr{syntactic position} and its kind its
\rr{sort}. Most of the time, these rules are invisible: they make
\code|x := x + 1| mean what it should. They become visible in error
messages, which mention them (``in a strong-unit''), and in the cases of
this chapter: the pointer comparisons, the order of uniting and widening,
and the annotations that help in an ambiguous situation.

\section{Lifetimes}\label{sec-lifetimes}

Every value that a program creates exists for a certain time. A cell
allocated by \code|loc|, which includes every variable declaration,
lives as long as the block in which it is allocated is running; when the
block ends, the cell is gone. A cell allocated by \code|heap| lives as
long as the program can reach it; when no reference to it remains, its
memory is reclaimed by a garbage collector. A function can therefore
return a reference to a new heap cell:

\example{ch07/heap}

The Report calls the lifetime of a value its \rr{scope}, a word this book
reserves for the visibility of names (Section~\ref{sec-scope}). A block
creates a lifetime that is newer, that is shorter, than the lifetime of
the enclosing block; heap cells and the values of the standard library
have the longest lifetime. The lifetime of a record or an array that
contains references is that of its shortest-lived reference.

The \term{lifetime rule} says that a reference must never outlive the
cell it refers to: a value may not be assigned to a variable that lives
longer than the value. Storing a reference to a local cell in a variable
of an enclosing block, or returning it from the block, breaks the rule:
\begin{lstlisting}
ref int p;
begin
   int x := 1;
   p := x          { wrong: p outlives x }
end
\end{lstlisting}
Rust enforces a rule of the same kind with its borrow checker, entirely
at compile time; in Algol 68 it can only partly be checked by the
compiler and otherwise must be checked while the program runs.

\begin{implnote}
The Report says that a program that breaks the lifetime rule is wrong and
that its further behavior is undefined, which allows but does not
require an implementation to check it. Algol 68 Genie checks: it warns at
compile time and stops the program with the error ``value is exported out
of its scope''. ga68 and a68j do not check, and such a program may even
appear to work, which makes the error all the more dangerous. Use
\code|heap| for anything that must outlive its block.
\end{implnote}

\subsection*{The lifetime of a \textbf{loc} cell}

\code|loc| is an expression like any other and can appear anywhere a
reference is needed, for example as a dummy argument for a
\code|ref| parameter whose result is not wanted. The precise rule for its
lifetime has a surprising consequence. A \code|loc| cell belongs to the
nearer of two things: the function or operator body in which it is
allocated (every call counts), or the innermost block, loop body or
other range that declares a variable, a constant, a function, an
operator or a type. Ranges without such
declarations do not count, nor do the loop counter, the name in an
alternative of a conformity clause, labels and priority declarations. Every call of
a function counts, even if its body declares nothing:
\code|proc f = ref int: loc int := 0| returns a reference to a cell that
is gone when the call is over. The following program keeps references to cells allocated in
a loop body:

\program{ch07/loop-loc}

\noindent Algol 68 Genie stops it:
\outputof{ch07/loop-loc}

The loop body declares \code|square|, so every \code|loc int| belongs to
one execution of the body and is gone when the next one starts. Without
the declaration, writing \code|cells[i] := loc int := i * i|, the cells
belong to the block around the loop, and the program is correct and
prints 4. A harmless-looking change to the body thus changes the
lifetime of the cells. The advice is simple: if a reference leaves the
block or loop body in which its cell is allocated, allocate the cell with
\code|heap|.

\code|heap| can also be used in declarations: \code|heap int n := 5| and
\code|heap [10] int a| declare variables whose cells live on the heap.
There is no operation to free a heap cell; that is left to the garbage
collector, which costs time and memory. The distinction is the one
between the stack and the heap in C++ and Rust (\texttt{Box}); Java and
Python allocate every object on the heap.

\section{Lifetimes of functions}\label{sec-proc-lifetimes}

The lifetime rule also applies to functions, and this is where Algol 68
differs most from current languages. A function value lives as long as
the innermost block that declares a name the function uses but does not
declare itself: a variable, a constant, a parameter of an enclosing
function, another function, an operator, a label it jumps to, or a type
name with bounds. A function that uses such a
name cannot be returned from that block or stored in a variable outside
it.

In JavaScript, Python or Haskell, returning a function that composes two
other functions is an everyday idiom. In Algol 68 it is an error:

\program{ch07/compose}

\noindent Algol 68 Genie stops at the call of \code|compose|:
\outputof{ch07/compose}

The inner function uses \code|f| and \code|g|, the parameters of
\code|compose|, which exist only while \code|compose| runs. Putting the
data on the heap does not help either:
\begin{lstlisting}
proc int next = (heap int count := 0; int: count +:= 1)
\end{lstlisting}
The cell of \code|count| would live long enough, but the name
\code|count| belongs to the block, and the function uses the name. Algol
68 Genie stops this program as well, with the same message; ga68 and a68j
run both programs and print plausible results. Here the lack of a check
is most dangerous, since such a program appears to work.

The reason for the rule is that an Algol 68 function does not keep alive
the variables and constants it uses, as a closure in Lisp, JavaScript or
Python does by keeping them on the heap; it refers to the block in which
they are declared, which lives on the stack. A C++ lambda that captures by reference has the same
problem, without any check; Rust rejects such a closure at compile time
unless it is declared \code|move|.

The usual way out is to pass the state explicitly. A counter becomes a
heap cell and a function that takes it as a parameter:

\example{ch07/counter}

\noindent which prints \code|1 2 3 1|. A function that composes two
functions becomes a procedure \code|(Fn f, g, real x) real: f (g (x))|
that receives all three values at each call. Where a function value must
be kept, a record that holds the state together with a function that
takes the state as a parameter does the job of a closure. Algol 68 Genie
and a68j offer a further way out, partial parametrization, an extension of
the language that Appendix~\ref{app-partial} describes.

\section{Recursive types}\label{sec-recursive}

A type declaration may refer to the type it declares, which allows
linked data structures. A list of integers is a chain of records, each of
which holds a value and a reference to the next record:

\example{ch07/list}

The constant \code|no_node| is \code|nil| of the type \code|ref Node|;
naming it makes the comparisons readable. The loop builds the list by
putting each new record in front: \code|heap Node := (i * i, list)|
allocates a record on the heap, initializes it, and yields the reference
to it, which is stored in \code|list|. The second loop follows the
references; \code|value of p| and \code|next of p| dereference \code|p|
automatically. The records must be allocated with \code|heap|: with
\code|loc| inside a procedure, the list would refer to cells that are
gone when the procedure returns, and ga68 and a68j would not report it.

The recursion must pass through a reference or a function type. A record
that contains itself directly would be infinitely large, and the compiler
rejects it:

\example{ch07/bad-mode}

Passing through a reference is not enough on its own, though:
\code|mode A = ref A| or \code|mode U = union (int, ref U)| would make a
type out of nothing but references, and they are rejected as well. The
precise rule is that every cycle in a type definition must pass both
through a reference or a function type, and through a record or a
function type with parameters. \code|struct (int value, ref Node next)|
satisfies it, and so does \code|mode F = proc (F) int|. The Report calls
the admissible recursive types \rr{well-formed modes}.

Several types can refer to each other, if they are declared in the same
block:

\example{ch07/mutual}

\subsection*{Pointers to links}

The type \code|ref ref Node| is the type of a variable that holds a
link: the variable \code|list| itself, or the field \code|next| of a
node. A procedure that receives such a reference can change the link,
which makes operations on lists simple. The following procedure removes
the first node with a given value. It walks along the links, not along
the nodes, so that the first node needs no special treatment:

\example{ch07/remove}

After \code|remove (list, 1)|, \code|remove (list, 3)|,
\code|remove (list, 5)| and \code|remove (list, 9)| on the list 1 to 5,
the list holds \code|2 4|. \code|link| is a pointer variable of type
\code|ref ref ref Node|. \code|ref Node (link)| is the node that the
current link refers to, and \code|next of link| is the link inside that
node. The assignment \code|ref ref Node (link) := next of link| changes
the current link so that it skips the node.

\begin{pitfall}
The annotation decides what is assigned. With
\code|ref Node (link) := next of link|, the left-hand side would be the
node itself, and the assignment would copy the next node's value and
link into it: a different operation, which fails at the last node, whose
\code|next| is \code|nil|. Both versions compile. When an assignment
involves several levels of references, it pays to write down the type of
each side.
\end{pitfall}

\section{Unions and recursive types}\label{sec-expr}

Unions and recursive types together describe tree-shaped data whose
nodes are of different kinds: the syntax tree of a compiler, a document,
a JSON value. An arithmetic expression is a number, a variable, or an
operator with two operands:

\program{ch07/expr}

\noindent prints
\outputof{ch07/expr}

The type \code|Expr| is a union of \code|real|, \code|char| (the name of
a variable) and \code|ref Binop|, and \code|Binop| is a record that
contains two \code|Expr| values. The cycle passes through the reference
and the record, so the types are well-formed. \code|union (real, char, Binop)|
without the reference would be infinitely large and is rejected. In Rust
the same type is written
\code|enum Expr { Num(f64), Var(char), Bin(Box<Binop>) }|; the
\texttt{Box} plays the role of \code|ref|.

\code|show| and \code|eval| are recursive functions with one alternative
of a conformity clause per kind of node. The procedure \code|bin|
allocates a node on the heap; its arguments are united to \code|Expr|
automatically, since an argument is in a strong context.

\subsection*{Operators on unions}

It is tempting to make the construction of expressions more readable
with operators: \code|op + = (Expr a, b) Expr: bin ("+", a, b)|, so that
an expression can be written as \code|3.0 * x * x - 2.0 * x + 1.0|.
Declared in the program above, this operator takes over every
\code|+| of two \code|real| values in its scope. The operands of an
operator are in a firm context, where uniting is allowed, and the new
\code|+| is declared in an inner block, so it hides the standard
\code|+| of the library for \code|real| operands. \code|l + r| in
\code|eval| then builds an \code|Expr| instead of adding, and the
compiler reports
\begin{lstlisting}[style=output]
error: 'EXPR'mode cannot be coerced to 'real' in a strong-unit
\end{lstlisting}
Inside the body of an operator on a union, the effect can be worse: an
\code|op = = (union (int, Rational) a, b) bool| whose body compares two
\code|int| values with \code|=| calls itself, without end. The remedy is
to declare an operator for each combination of member types that is
needed (\code|(Rational, Rational)|, \code|(int, Rational)|,
\code|(Rational, int)|), or to give the operator a new name, such as
\code|PLUS|.

\section{When two types are the same}\label{sec-equivalence}

Section~\ref{sec-modes} showed that type names are only abbreviations:
two record types with the same fields are the same type. This also holds
for recursive types. The types
\begin{lstlisting}
mode A = struct (int v, ref A next);
mode B = struct (int v, ref B next)
\end{lstlisting}
are the same type, since expanding either definition gives the same
infinite structure. Deciding this is a non-trivial task for the compiler,
but it is invisible to the programmer. Together with the rules for unions,
it means that two types are the same if and only if they have the same
structure; the Report calls them \rr{equivalent modes}.

\section{A complete program}\label{sec-advanced-complete}

The final program counts the words of a text with a binary search tree.
Each node holds a word, a count and two subtrees. The procedure
\code|insert| takes a reference to the variable that holds a subtree, of
type \code|ref ref Tree|, so that it can replace an empty subtree with a
new node; \code|show| traverses the tree in order and prints the words
alphabetically.

\example{ch07/tree}

The test \code|ref Tree (t) is no_tree| uses the cast of
Section~\ref{sec-pointers}: \code|t| is a reference to a variable, and the
question is whether the variable holds \code|nil|. In \code|show|, the
parameter is the tree itself, so \code|t isnt no_tree| needs no cast.
The splitting loop uses \code|orel| (Section~\ref{sec-logic}) so that
\code|text[i]| is not evaluated when \code|i| is beyond the end.

\section{A larger example: a JSON parser}\label{sec-json}

The programs so far have been short, each one showing a few features.
This section combines the features of this and the earlier chapters in a
program of about 230 lines. It has a parser that reads JSON text into a
tree of values, a printer that writes the tree back with indentation,
and an operator that selects parts of the tree. Section~\ref{sec-expr}
mentioned JSON, the data format of most web services, as an example of
tree-shaped data whose nodes are of different kinds. Algol 68 can
describe such data directly.

\subsection*{The types}

\programpart{ch07/json}{types}

A JSON value is of one of six kinds, and \code|Json| is a union of six
types. \code|null| is the value \code|empty| of type \code|void|
(Section~\ref{sec-union-details}). \code|true| and \code|false| are
\code|bool| values, and numbers are \code|real|. Arrays and objects are
references to records on the heap. An object is an array of members,
each a key with a value, kept in the order of the text. \code|Error|
describes a syntax error found by the parser.

The types are recursive: a \code|Json| value can be a reference to an
\code|Array|, which contains \code|Json| values. Every cycle passes
through a reference and through a record, so the types are well-formed
(Section~\ref{sec-recursive}). The shorter \code|union (..., ref [] Json)|
would be rejected. Its cycle passes through a reference but through no
record, like the cycle of \code|union (int, ref U)|. So the record with
its single field is needed to make the type well-formed. Since arrays and
objects are held by reference, copying a \code|Json| value copies only
the reference, as with the objects of Java or Python. Since the field is
a resizable array, a record on the heap can grow while the parser fills
it. In Rust the type would be written
\begin{lstlisting}[language={}]
enum Json { Null, Bool(bool), Number(f64), String(String),
            Array(Vec<Json>), Object(Vec<(String, Json)>) }
\end{lstlisting}
\noindent where \texttt{Vec} is a resizable array on the heap.

\subsection*{Growing an array}

\programpart{ch07/json}{append}

The operator \code|+:=| appends a value to a resizable array, like the
function \code|append| of Section~\ref{sec-flex}. As an operator it
reads like the \code|s +:= c| of strings: \code|elements of a +:= value|.
It is declared twice, once for each element type, because Algol 68 has
no generic functions. The standard \code|+:=| for strings remains
available, since an operator is chosen by the types of its operands
(Section~\ref{sec-op-identification}). Copying the array for every new
element takes quadratic time in the number of elements. For large
documents, the doubling of Section~\ref{sec-vector} would avoid that.

\subsection*{The parser}

The parser works by \term{recursive descent}: there is one function for
each kind of value, and it calls the functions for the parts of the
value. The functions are declared inside \code|parse| and share the text
and the current position \code|pos|, as the methods of an object share
its fields:

\programpart{ch07/json}{parse}

\code|peek| skips white space and yields the next character, or
\code|REPR 0| at the end of the text. \code|expect| checks one character
and steps over it. The condition of the \code|if| in \code|value|
declares \code|c|, which is visible in all branches
(Section~\ref{sec-if}). The branches yield values of different types: a
reference to an object or an array, a string, a number, a \code|bool|
or \code|empty|. The \code|if| as a whole has the type \code|Json|, the
result type of \code|value|, and each branch is united to it. The last
branch calls \code|fail|, which never returns. Its type is \code|void|,
so \code|skip| supplies a \code|Json| value; that value is never used.

\programpart{ch07/json}{containers}

\code|object| and \code|array| create a record on the heap and append to
its field while they read the elements. \code|value| calls \code|array|
and \code|array| calls \code|value|. The functions are mutually
recursive, and the depth of the recursion follows the nesting of the
text. \code|object| and \code|array| are declared after \code|value|,
which uses them. That works because a function declared with a function
literal exists from the start of its block (Section~\ref{sec-scope}).

\code|member| reads the key in a declaration of its own and checks the
colon before the display \code|(key, value)| calls \code|value|. The
parts of a display are evaluated in no particular order
(Section~\ref{sec-collateral}). A display \code|(string_value, value)|
might therefore read the value before the key.

\programpart{ch07/json}{scalars}

\code|string_value| collects the characters up to the closing quotation
mark. It turns the escapes \texttt{\textbackslash n} and
\texttt{\textbackslash t} into a newline and a tab. Any other character
after a backslash stands for itself, which covers \texttt{\textbackslash"}
and \texttt{\textbackslash\textbackslash}. In an Algol 68 string, a
quotation mark is written twice, so \code|""""| is a single quotation
mark; the backslash has no special meaning.

\code|number| checks the syntax of a JSON number: an optional minus
sign, digits, optionally a point and more digits, and optionally an
exponent. Each part is one line of code, so the function reads like the
grammar. \code|next_is| uses the library function
\code|char_in_string| (Appendix~\ref{app-library}). \code|loc int|
serves as a dummy argument for the position, which is not needed
(Section~\ref{sec-lifetimes}). The transput converts the digits into a
\code|real|: \code|associate| turns the string into a file
(Section~\ref{sec-associate}), and \code|get| reads the number from it.

\programpart{ch07/json}{parse-end}

The body of \code|parse| reads one value and checks that nothing but
white space follows it. If all goes well, \code|exit| ends the body with
\code|doc| (Section~\ref{sec-jumps}). If anything goes wrong,
\code|fail| records the position and the message and jumps to the label
\code|failed|. The jump leaves all active calls of \code|value|,
\code|array|, \code|object| and the other functions at once, however
deep the recursion is, and \code|parse| yields the error. Other languages
use exceptions for this. In Algol 68 it is an ordinary jump to a label in
an enclosing block, and \code|fail| can jump there because it is declared
in that block. The result type \code|union (Json, Error)| is a union of
seven types, since a union that contains a union contains its members
(Section~\ref{sec-unions}).

\subsection*{Printing}

\programpart{ch07/json}{show}

\code|show| is a conformity clause with one alternative for each kind of
value; \code|(void)| is chosen for \code|null|. For arrays and objects,
\code|show| calls itself for the elements with a deeper indentation. An
array whose elements are all scalars is printed on one line. The test is
\code|scalar|, a conformity clause without names. Two small functions
turn numbers and strings into text:

\programpart{ch07/json}{text}

\code|number_text| writes an integral number without a fraction. For
other numbers it drops the trailing zeros of \code|fixed|. JSON requires
a digit before the point, which \code|fixed| leaves out:
\code|fixed (0.5, 0, 6)| yields \code|.500000|. \code|quote| writes the
escapes back.

\subsection*{Selecting parts of a value}

\programpart{ch07/json}{select}

\code|SEL| selects a member of an object by its key or an element of an
array by its index. It yields \code|null| if there is none, so that
selections can be chained without checks:
\code|doc SEL "implementations" SEL 3 SEL "name"|. Operators of the same
precedence are applied from left to right (Section~\ref{sec-precedence}),
so the chain is read as
\code|((doc SEL "implementations") SEL 3) SEL "name"|. The operator has
two declarations, one with a \code|string| and one with an \code|int|
as the right operand.

\code|SEL| needs a name of its own. An operator \code|/| with the
operands \code|Json| and \code|int| would also be chosen for \code|x / 2|
with a \code|real| \code|x|, because the \code|real| operand is united
to \code|Json| (Section~\ref{sec-expr}). A key of one letter must be
written \code|string ("a")|: \code|"a"| alone is a \code|char|, and no
\code|SEL| takes a \code|char| (Section~\ref{sec-union-details}).

\subsection*{The main program}

\programpart{ch07/json}{main}

\noindent prints
\outputof{ch07/json}

Since a quotation mark is written twice in a string literal, the JSON
text contains many of them. The conformity clause on the result of
\code|parse| has an alternative for a union, \code|(Json doc)|, which is
chosen for each of the six kinds of values
(Section~\ref{sec-union-details}). The number \code|-2.5e1| is printed
as \code|-25|, and error positions count characters from 1.

\subsection*{What the example shows}

The parser combines several features:
\begin{itemize}
\item tagged unions, with conformity clauses that check the type at run
  time;
\item recursive types on a heap whose memory is managed automatically;
\item nested functions that share the variables of their block;
\item operators with user-defined precedence;
\item a jump out of a deep recursion.
\end{itemize}
All of them belong to the language of the Revised Report of 1975. The
program goes beyond it only in the short-circuit \code|andth| and in the
transput of \code|Genie_Transput|. In the widely used languages, most of these
features came decades later: \texttt{std::variant} with C++17, the
\texttt{match} statement with Python~3.10 in 2021, and pattern matching
over sealed types with Java~21 in 2023. A C version of the program would
need a hand-written tag in every value and explicit memory management.

The program also shows the limits of the language. Algol 68 has no
generic functions, hence the two \code|+:=| operators, and it has no
exceptions, only jumps. A further limit comes from the implementations:
with ga68 and a68j, \code|real| has single precision
(Section~\ref{sec-types}), so a number with more than about seven
significant digits loses its last digits. A parser for real use would
keep \code|long real| values.

\section*{Exercises}
\addcontentsline{toc}{section}{Exercises}

The \hyperref[ans-ch7]{answers} are at the end of the book.

\begin{enumerate}
\item What are the bounds of \code|s| after
  \code|flex [1:0] char s; s := "abc"; s := s[2:3] + "xyz"|?
\item Why can \code|ref [] int| not be the type of the parameter of the
  procedure \code|append|?
\item Write a procedure that takes a \code|union (int, real, bool)| and
  returns it as a \code|real|: an \code|int| widened, a \code|bool| as
  1.0 or 0.0.
\item Given \code|int x := 1; ref int p := x|, which of the following
  change \code|x|? (a) \code|p := 2| (b) \code|ref int (p) := 2|
  (c) \code|x := 2|
\item Why is \code|q is nil| always false for a variable \code|q| of
  type \code|ref int|, and what is the correct test?
\item In which contexts are the following expressions:
  (a) \code|x| in \code|x := 1|, (b) \code|x| in \code|y := x|,
  (c) \code|x| in \code|x + 1|, (d) \code|b| in \code|if b then|?
\item Write a procedure that reverses a list of the type \code|Node| of
  Section~\ref{sec-recursive} in place.
\item \code|flex [1:0] string words| holds some words. Why is
  \code|words[1] +:= "s"| rejected, and how can the program add an
  \code|s| to the first word?
\item Which of these declarations are accepted?
  (a) \code|union (int, string) u := "a"|
  (b) \code|union (int, string) u := "ab"|
  (c) \code|union (real, string) u := 2|
  (d) \code|union (int, real) u := 2|
\item With \code|int a := 1, b := 1; ref int p := a, q := b|, what do
  \code|p is q|, \code|ref int (p) is ref int (q)| and
  \code|ref int (p) is a| yield? What changes after \code|q := a|?
\item The \code|push| of Section~\ref{sec-vector} doubles the capacity.
  How many elements are copied in total by 1000 calls of \code|push| on
  an empty \code|Vector|, and how many by 1000 calls of \code|append|?
\item Why is the following procedure wrong, and how do ga68, a68j and
  Algol 68 Genie behave when it is used?
\begin{lstlisting}
proc adder = (int n) proc (int) int: (int x) int: x + n
\end{lstlisting}
\item Does the program of Section~\ref{sec-lifetimes} that stores
  \code|loc| cells remain correct if the body of the loop is
  \code|cells[i] := loc int := i * i; print (i)|? And if it is
  \code|int j = i; cells[i] := loc int := j|?
\item Build two lists of type \code|ref Node| with the same contents
  1, 2, 3. What does \code|a is b| yield? Write an operator \code|EQ|
  that compares two lists element by element.
\end{enumerate}

\section*{Programming problems}
\addcontentsline{toc}{section}{Programming problems}

\begin{enumerate}
\item \textbf{Text in pieces.} Keep a long text as a linked list of
  pieces, each a \code|string|, so that joining two texts only links
  their lists. Write procedures for joining, for the length, for the
  character at a given position, and for turning the list back into one
  \code|string|. Which operations are cheaper than with a single
  \code|string|, and which are more expensive?
\item \textbf{List operations.} For the lists of
  Section~\ref{sec-recursive}, write procedures that return the
  $n$-th element, append a copy of one list to another, merge two sorted
  lists into one, and sort a list by repeatedly merging.
\item \textbf{Sets.} Represent sets of integers as sorted linked lists,
  and write the operators \code|+| (union), \code|*| (intersection),
  \code|-| (difference) and \code|ELEM| (membership, as in
  \code|3 ELEM s|). Then write the same operators for sets represented
  as \code|bits| (Section~\ref{sec-bits}) and compare.
\item \textbf{Tree sort.} Sort an array of integers by inserting its
  elements into a binary search tree and reading them back in order.
  Duplicates must be kept. What is the cost for an array that is
  already sorted?
\item \textbf{Postfix expressions.} Read an expression in postfix
  notation (\texttt{3 4 + 2 *}) and evaluate it with a stack of
  \code|real| values implemented as a linked list.
\item \textbf{Parsing.} Extend the program of Section~\ref{sec-expr}:
  read a fully parenthesized expression such as
  \texttt{((3 * (x * x)) - (2 * x))}, build its tree, print it, and
  evaluate it for several values of \code|x|.
\item \textbf{Symbolic differentiation.} Write a procedure that returns
  the derivative of an \code|Expr| with respect to \code|x|, and a
  procedure that simplifies the result: $0 \cdot e = 0$, $1 \cdot e = e$,
  $e + 0 = e$, and operations on two numbers are computed.
\item \textbf{Optional results.} Write a procedure that searches a
  list of (name, number) records for a name and returns
  \code|union (int, void)|. Use it in a program that reads names and
  reports each one's number or that it is unknown, with an \code|out|
  branch in every conformity clause.
\item \textbf{A hash table.} Implement a table from strings to integers
  as an array \code|[0 : m - 1] ref Entry| of linked lists, with
  procedures to insert, look up and delete. Count the words of a text
  with it and compare the result with the tree of
  Section~\ref{sec-advanced-complete}.
\item \textbf{Pascal's triangle.} Store the first 15 rows of Pascal's
  triangle in a jagged array, once in a \code|flex [1:0] flex [1:0] int|
  that grows row by row, and once as a \code|[15] ref [] int| whose rows
  are allocated with \code|heap|.
\item \textbf{A deck of cards.} Declare the 52 cards of a deck as a
  fixed array of records, each with a field \code|ref Card next|, and
  link them into a circle. A hand is a record with a count and a
  reference to its first card. Deal cards from the deck to four hands
  only by changing links, without copying a card. Why is it allowed here
  to keep references to elements of the array?
\item \textbf{A queue.} Implement a queue as a record with references
  to its first and its last node, so that both adding at the end and
  removing at the front take constant time. Write a procedure that
  creates a candidate node on the heap and appends it only if it passes a
  test; what happens to a rejected node?
\item \textbf{Exact or approximate.} Find the largest $n$ whose
  factorial fits in an \code|int|, without causing an overflow. Then
  write a function that returns \code|union (int, real)|: the exact
  factorial up to that $n$, and Stirling's approximation beyond it. How
  must a caller use the result?
\item \textbf{More JSON.} Extend the parser of
  Section~\ref{sec-json}. (a) Report errors with line and column instead
  of the position. (b) Accept the escape \texttt{\textbackslash u}
  followed by four hexadecimal digits, using \code|REPR|. This works
  with ga68 and a68j; in Algol 68 Genie, \code|max_abs_char| is 255.
  (c) Write an operator \code|EQ| that compares two JSON values, where the
  members of two objects may be in a different order. (d) Add \code|int|
  to the union, for numbers without a point and an exponent. Which
  conformity clauses must change, and does the compiler point them out?
\item \textbf{Cross-reference.} Read an Algol 68 program and build an
  index of its identifiers: a tree, sorted by name, whose nodes hold a
  list of the line numbers where the name occurs. Print the index in
  alphabetical order.
  Reserved words, which are written in lower case in SUPPER stropping,
  can be recognized with a table.
\end{enumerate}
